% 天王星（综述）
% license CCBYSA3
% type Wiki

本文根据 CC-BY-SA 协议转载翻译自维基百科\href{https://en.wikipedia.org/wiki/Uranus}{相关文章}。
\begin{figure}[ht]
\centering
\includegraphics[width=8cm]{./figures/fb24dd1c06777ad6.png}
\caption{由旅行者 2 号拍摄的天王星真彩图像\(^\text{[a]}\)。其苍白而柔和的外观源于覆盖在云层之上的一层薄雾。} \label{fig_Uranus_1}
\end{figure}
天王星是距离太阳的第七颗行星。它是一颗呈青色的气态冰巨行星。该行星的大部分由处于超临界相状态的水、氨和甲烷构成，天文学上将这些物质称为“冰”或挥发物。该行星的大气具有复杂的分层云结构，并在所有太阳系行星中具有最低的最低温度（49 K（−224 °C; −371 °F））。它具有 82.23° 的显著自转轴倾角，自转周期为 17 小时 14 分钟，且为逆行自转。这意味着在其围绕太阳完成一个为期 84 个地球年的公转周期期间，其两极会经历约 42 年的连续日照，随后是约 42 年的连续黑暗。

在太阳系行星中，天王星直径排名第三、质量排名第四。根据当前模型，在其挥发性地幔层内部存在一个岩质核心，外围包覆着厚厚的氢和氦大气。在其高层大气中检测到痕量的碳氢化合物（被认为通过水解作用产生）以及一氧化碳和二氧化碳（被认为源自彗星）。天王星大气中存在许多尚未解释的气候现象，例如其峰值风速 900 km/h（560 mph）\(^\text{[24]}\)、极冠的变化以及其不规则的云层结构。该行星的内部热量相较其他巨行星异常低，其原因尚不清楚。

与其他巨行星一样，天王星拥有环系统、磁层以及众多天然卫星。其极为暗淡的环系统仅反射约 2\% 的入射光。天王星的 29 颗天然卫星中包含 19 颗已知的规则卫星，其中 14 颗为小型内侧卫星。在更外侧的是该行星五颗较大的主要卫星：米兰达（Miranda）、艾丽尔（Ariel）、安布里尔（Umbriel）、泰坦妮亚（Titania）和奥伯龙（Oberon）。在距离天王星更远的轨道上则有十颗已知的不规则卫星。该行星的磁层高度不对称，且含有大量带电粒子，这可能是其环和卫星变暗的原因。

天王星肉眼可见，但亮度极低，并且相对于背景恒星移动非常缓慢，因此直到 1781 年被威廉·赫歇尔首次观测到之前，它从未被归类为一颗行星。在其发现约七十年后，人们达成共识，将其命名为希腊原始神祇之一乌拉诺斯（Uranus，Ouranos）。截至 2025 年，它仅被探测器访问过一次，即 1986 年旅行者二号（Voyager 2）飞掠该行星\(^\text{[25]}\)。尽管如今可通过望远镜分辨并观测其细节，但科研界强烈希望再次造访该行星——这一点体现于《行星科学十年调查》将拟议的“天王星轨道器与探测器任务”列为 2023–2032 年的最高优先级任务之一，以及中国国家航天局提出利用天问四号的子探测器飞掠天王星的计划\(^\text{[26]}\)。
\subsection{历史}
1781 年 3 月 13 日（天王星被发现之日），天王星的位置（以叉号标示）与古典行星一样，天王星肉眼可见，但由于其亮度极低且公转缓慢，古代观测者从未将其识别为一颗行星\(^\text{[27]}\)。威廉·赫歇尔于 1781 年 3 月 13 日首次观测到天王星，由此发现了这颗行星，这也是人类历史上首次扩展太阳系已知边界，并使天王星成为第一颗借助望远镜而被确认的行星。天王星的发现实际上使当时已知太阳系的尺度扩大了一倍，因为天王星距离太阳的平均距离约为土星的两倍。
\subsubsection{发现}
在天王星被确认是一颗行星之前，它曾被多次观测到，但一般被误认为是一颗恒星。显然，喜帕恰斯在公元前 128 年观测过它，当时他为自己的星表测量恒星的位置，该星表后来被纳入托勒密的《天文大成》中。该星表给出了处于室女座的四颗恒星的位置，它们构成一个四边形。其中有一颗恒星实际上并不存在，而在公元前 128 年 4 月，天王星正位于该位置\(^\text{[28]}\)。最早的确定无误的观测记录是在 1690 年，当时约翰·弗拉姆斯蒂德至少六次观测到它，并将其编入星表为金牛座 34 号星。詹姆斯·布拉德雷在 1748 年、1750 年和 1753 年三次观测到它。托比亚斯·迈耶于 1756 年观测到它一次。法国天文学家皮埃尔·夏尔·勒莫尼耶在 1750 年至 1769 年期间至少十二次观测到天王星，其中包括连续四个夜晚的观测\(^\text{[29]}\)。

威廉·赫歇尔于 1781 年 3 月 13 日，在英国萨默塞特郡巴斯新国王街 19 号自家花园（现为赫歇尔天文学博物馆）观测到了天王星\(^\text{[30]}\)，并于 1781 年 4 月 26 日最初将其报告为一颗彗星\(^\text{[31]}\)。使用一架自制的 6.2 英寸反射望远镜，赫歇尔“从事了一系列关于恒星视差的观测”\(^\text{[32][33]}\)。

赫歇尔在日记中写道：“在靠近金牛座 ζ 的四分仪区域……要么是一颗星云状恒星，要么可能是一颗彗星”\(^\text{[34]}\)。3 月 17 日他记道：“我寻找那颗彗星或星云状恒星，并发现它是一颗彗星，因为它改变了位置”\(^\text{[35]}\)。当他向英国皇家学会呈报他的发现时，他依然坚持认为自己发现的是一颗彗星，但也含蓄地将其与一颗行星作比较\(^\text{[32]}\)：

“当我第一次看到这颗彗星时所使用的放大倍率是 227。根据经验我知道，恒星的视直径在更高倍率下不会按比例放大，而行星则会；因此我随即将倍率调至 460 和 932，发现彗星的直径随倍率按比例增大，这正符合其并非恒星的假设，而与之比较的那些恒星的直径并未按相同比例增大。此外，这颗彗星在远超其亮度所能承受的高倍率下，显得模糊且轮廓不清，而恒星则保持了我从成千上万次观测中所熟知的那种光泽与清晰度。随后的一切证明了我的猜测是正确的，这正是我们最近观测到的那颗彗星。”

赫歇尔将他的发现通知了皇家天文学家内维尔·马斯克林，并在 1781 年 4 月 23 日收到他困惑的回信：“我不知道该如何称呼它。它既有可能是一颗沿着近圆形轨道绕太阳运行的常规行星，也有可能是一颗沿着极端椭圆轨道运行的彗星。我尚未看到它有任何彗发或彗尾”\(^\text{[36]}\)。

尽管赫歇尔继续将他的新天体描述为一颗彗星，但其他天文学家已经开始怀疑并非如此。芬兰—瑞典籍、在俄罗斯工作的天文学家安德斯·约翰·勒克塞尔是第一位计算出这一新天体轨道的人\(^\text{[37]}\)。其近乎圆形的轨道暗示它是一颗行星而非一颗彗星。柏林天文学家约翰·埃勒特·博德将赫歇尔的发现描述为“一颗移动的恒星，可被视为迄今未知的、在土星轨道之外运行的类行星天体”\(^\text{[38]}\)。博德得出结论认为，其近圆形轨道与行星轨道更为相似，而非彗星轨道\(^\text{[39]}\)。

该天体很快被接受为一颗新行星。到 1783 年，赫歇尔向皇家学会会长约瑟夫·班克斯承认了这一点：“据欧洲最杰出天文学家的观测结果显示，那颗我在 1781 年 3 月有幸向他们指出的新星，是我们太阳系中的一颗主行星”\(^\text{[40]}\)。为表彰他的成就，乔治三世国王授予赫歇尔每年 200 英镑的津贴（相当于 2023 年的 30,000 英镑）\(^\text{[41]}\)，条件是他搬到温莎，以便王室成员能够通过他的望远镜观测天体\(^\text{[42]}\)。
\begin{figure}[ht]
\centering
\includegraphics[width=8cm]{./figures/221d0cedb3ff1f58.png}
\caption{} \label{fig_Uranus_3}
\end{figure}
\subsubsection{姓名}
\begin{figure}[ht]
\centering
\includegraphics[width=6cm]{./figures/57b83b23d4c37bca.png}
\caption{威廉·赫歇尔，天王星的发现者} \label{fig_Uranus_2}
\end{figure}
“天王星”这一名称源自古希腊的天空神乌拉诺斯（古希腊语：Οὐρανός），在罗马神话中称为 Caelus，为克洛诺斯（即土星）的父亲、宙斯（即木星）的祖父以及阿瑞斯（即火星）的曾祖父，其名称在拉丁语中被写作 Uranus（IPA: [ˈuːranʊs]）\(^\text{[2]}\)。它是八大行星中唯一一个英文名称来自希腊神话人物的行星。天文学家所偏好的 “Uranus” 发音是 /ˈjʊərənəs/（YOOR-ə-nəs）\(^\text{[1]}\)，其发音方式包含英语的长 “u”，且重音落在第一音节，与拉丁语 Uranus 相同；另一种常见发音 /jʊˈreɪnəs/（yoo-RAY-nəs）则将重音置于第二音节并使用长 a，两者皆被视为可接受\(^\text{[g]}\)。

关于行星名称的共识直到发现后的近 70 年后才最终达成。在最初的讨论中，马斯基林曾请求赫歇尔“请施惠于天文学界，为你亲自发现、令我们深表感激的这颗行星赐名”\(^\text{[44]}\)。作为回应，赫歇尔决定将其命名为 Georgium Sidus（乔治之星），或称“乔治行星”，以纪念他的新庇护人乔治三世国王\(^\text{[45]}\)。他在写给约瑟夫·班克斯的信中解释了这一决定\(^\text{[40]}\)：

在古代神话时代，水星、金星、火星、木星与土星被命名为行星，是以其主要英雄与神祇之名相称。在当今更具哲学精神的时代，再以同样方式为我们新的天体命名为朱诺、帕拉斯、阿波罗或密涅瓦，则未免不太合宜。任何特殊事件或显著事物，其首要考虑应为年代记载：若在未来时代，人们询问这颗最新发现的行星是何时被发现的？那么，“在乔治三世统治时期”将是一个令人十分满意的答案。

赫歇尔提出的名称在不列颠与汉诺威之外并不受欢迎，于是很快出现了替代方案。天文学家热罗姆·拉朗德建议将其命名为 Herschel 以纪念发现者\(^\text{[46]}\)。瑞典天文学家 Erik Prosperin 提出了 Astraea、Cybele（后来成为小行星名称）以及 Neptune（后来成为下一颗被发现的行星的名称）。来自哥廷根的格奥尔格·李希滕贝格也支持 Astraea（作 Austräa），但该女神传统上与处女座而非金牛座相关。其他天文学家支持 Neptune，希望以此纪念英国皇家海军在美国独立战争期间的胜利，他们提出将新行星命名为 Neptune George III 或 Neptune Great Britain；这一折衷方案也由莱克塞尔所建议\(^\text{[37][47]}\)。丹尼尔·伯努利则提出了 Hypercronius 与 Transaturnis。Minerva 亦曾被提议为名称\(^\text{[47]}\)。
\begin{figure}[ht]
\centering
\includegraphics[width=6cm]{./figures/997dfb34e98d04d5.png}
\caption{约翰·埃勒特·波德，这位提出“天王星”名称的天文学家} \label{fig_Uranus_4}
\end{figure}
在 1782 年 3 月的一篇论文中，约翰·埃勒特·波德提出了“Uranus”这一名称，即希腊天空之神 Ouranos 的拉丁化形式\(^\text{[48]}\)。波德认为，行星名称应遵循神话体系，以免与其他行星显得不同；而 Uranus 作为第一代泰坦之父，是一个合适的名字\(^\text{[48]}\)。他还指出，这个名字十分优雅，因为正如土星（Saturn）是木星之父，新行星也应以土星之父命名\(^\text{[42][48][49][50]}\)。然而，他似乎并不知道 Uranus 只是该神名的拉丁化形式，其真正的罗马对应神是 Caelus。1789 年，波德在皇家科学院的同事马丁·克拉普罗特为其新发现的元素命名为“铀”（uranium），以支持波德的命名选择\(^\text{[51]}\)。最终，波德的建议成为最广泛使用的名称，并在 1850 年成为通用称呼，当时最后坚持使用 Georgium Sidus 的英国航海年鉴局也改用 Uranus\(^\text{[49]}\)。

天王星有两个天文学符号。其一为最早提出者：⛢\(^\text{[h]}\)，由约翰·戈特弗里德·科勒于 1782 年在波德的请求下提出\(^\text{[52]}\)。科勒建议使用铂金的符号，而铂金仅在 30 年前才被科学描述；由于没有传统的炼金术符号，他提出使用 ⛢ 或 ⛢，即由行星–金属符号 ☉（金）与 ♂（铁）组合而来，因为铂（金之“白金”）常与铁共生。波德认为符号直立的版本 ⛢ 更符合其他行星的符号体系，同时依然保持区别性\(^\text{[52]}\)。在现代天文学中，该符号在（罕见的）使用符号的情况下最为常见\(^\text{[53][54]}\)。第二个符号 ♅\(^\text{[i]}\) 是拉朗德于 1784 年提出的；在写给赫歇尔的信中，拉朗德将其描述为“un globe surmonté par la première lettre de votre nom”（“一个上方加上您姓氏首字母的圆球”）\(^\text{[46]}\)。这一符号在占星学中几乎是通用的。

在英语流行文化中，由于 Uranus 的常见发音类似于“your anus”，由此衍生出大量幽默内容\(^\text{[55]}\)。

在其他语言中，天王星有多种名称。在中文（天王星）、日文（天王星 Tennōsei）、韩文（천왕성 Cheonwangseong）与越南语（sao Thiên Vương）中，其名称都直译为“天之王星”\(^\text{[56][57][58][59]}\)。在泰语中，其正式名称为 Dao Yurenat（ดาวยูเรนัส），与英语一致；其另一名称 Dao Maruettayu（ดาวมฤตยู，“Mṛtyu 星”）源于梵语中表示“死亡”的词 mrtyu（मृत्यु）。在蒙古语中，其名称是 Tengeriin Van（Тэнгэрийн ван），意为“天空之王”，反映了其同名神祇作为天界统治者的地位。在夏威夷语中，其名称为 Heleʻekela，即“Herschel”名称的夏威夷语音译\(^\text{[60]}\)。
\subsection{形成}  
一般认为，冰巨行星与气巨行星之间的差异源于它们的形成历史\(^\text{[61][62][63]}\)。太阳系被假设形成于一个被称为原太阳星云的旋转气尘盘中。星云中的大部分气体——主要由氢和氦构成——形成了太阳，而尘埃颗粒聚集在一起形成最初的原行星。随着这些行星体逐渐增长，其中一些最终吸积了足够的物质，使得它们的引力能够束缚住星云中残余的气体\(^\text{[61][62][64]}\)。它们吸附的气体越多，体积就越大；体积越大，又能吸附更多气体，直到达到一个临界点，其大小开始呈指数式增加\(^\text{[65]}\)。冰巨行星仅吸附了数个地球质量的星云气体，因而从未达到这一临界点\(^\text{[61][62][66]}\)。近期关于行星迁移的模拟表明，两颗冰巨行星形成时的位置比现在更靠近太阳，并在形成后向外迁移（即 Nice 模型）\(^\text{[61]}\)。
\subsection{轨道与自转} 
天王星绕太阳一周需要 84 年。自 1781 年被发现以来\(^\text{[67]}\)，依据其相对于背景恒星的位置计算，该行星已经两次回到其最初被发现的位置——在金牛座双星 ζ Tauri 东北方向——分别是在 1865 年 3 月与 1949 年 3 月，并将在 2033 年 4 月再次回到该位置\(^\text{[68]}\)。
\begin{figure}[ht]
\centering
\includegraphics[width=8cm]{./figures/0658c53f45263d76.png}
\caption{在相隔 84 年的两个各为 2 年区间中，天王星的赤经变化} \label{fig_Uranus_5}
\end{figure}
由于其周期非常接近 84 年，因此它在恒星背景中的视位置与整整 84 年前几乎相同（见图）。天王星的逆行运动使其每年都会回到大致与前一年最偏东的位置相近的地方。

它与太阳的平均距离约为 20 AU（30 亿千米；20 亿英里）。其近日点与远日点之间的差值为 1.8 AU，大于任何一颗行星，但仍小于矮行星冥王星的对应差值\(^\text{[69]}\)。太阳光强度与距离的平方成反比——在天王星处（距离约为地球至太阳距离的 20 倍），光照强度约为地球的 1/400\(^\text{[70]}\)。

天王星的轨道要素最早由皮埃尔–西蒙·拉普拉斯于 1783 年计算\(^\text{[71]}\)。随着时间推移，预测轨道与观测轨道之间开始出现偏差，1841 年，约翰·考奇·亚当斯首次提出这些差异可能来自某颗未被观测到的行星的引力扰动。1845 年，乌尔班·勒威耶展开了对天王星轨道的独立研究。1846 年 9 月 23 日，约翰·戈特弗里德·加勒在几乎等同于勒威耶预测的位置发现了一颗新行星，即后来命名的海王星\(^\text{[72]}\)。

天王星内部的自转周期为 17 小时 14 分钟 52 秒\(^\text{[14]}\)，这一数值通过追踪天王星极光的自转运动而得出\(^\text{[73]}\)。与所有巨行星一样，它的高层大气在自转方向上出现强风。在某些纬度（例如大约南纬 60 度），可见大气特征移动得更快，仅需约 14 小时即可完成一次完整自转\(^\text{[74]}\)。
\subsubsection{自转轴倾角}
\begin{figure}[ht]
\centering
\includegraphics[width=8cm]{./figures/a3f72c0a3302a344.png}
\caption{从 1986 年到 2030 年，模拟从地球看到的天王星视图：从 1986 年的南半球夏至，经 2007 年的昼夜平分点，到 2028 年的北半球夏至。} \label{fig_Uranus_6}
\end{figure}
天王星的自转轴大致与太阳系的平面平行，其自转轴倾角可描述为 82.23° 或 97.77°，取决于哪个极被视为北极\({^\text{[j]}}\)。前一种定义遵循国际天文学联合会的规定，即北极是位于太阳系不变平面地球北侧的那一极。在这种定义下，天王星具有逆行自转。或者，若按照另一种惯例，以“旋转方向的右手定则”来定义天体的南北两极，则天王星的自转轴倾角可表示为 97.77°，这将反转哪个极被视为北极、哪个极被视为南极，并使得该行星呈现顺行自转\({^\text{[75]}}\)。这一特性使其季节变化完全不同于其他行星。冥王星与小行星 2 号智神星也具有极端的自转轴倾角。接近至日时，其中一个极持续朝向太阳，另一个极持续背向太阳，只有赤道附近一条狭窄的区域经历快速的昼夜交替，而太阳在地平线上非常低。随着天王星沿轨道移动至另一侧，两极朝向太阳的状态将完全反转。每个极大约经历 42 年的连续白昼，随后是 42 年的连续黑夜\({^\text{[76]}}\)。接近春分或秋分时，太阳照射天王星赤道，使其出现类似其他行星所经历的昼夜交替。

这种自转轴取向的一个结果是：在一个天王星年内，天王星近极地区接收到的太阳能量平均多于赤道地区。然而，天王星的赤道区域却比极地区更热。造成这一现象的内在机制尚不清楚。天王星异常自转轴倾斜的原因也并未确定，但通常的推测是：在太阳系形成期间，一个具有地球大小的原行星与天王星发生碰撞，从而导致了这种偏斜的取向\({^\text{[77]}}\)。达勒姆大学 Jacob Kegerreis 的研究表明，这一倾斜可能是由于一个比地球更大的原行星在 30–40 亿年前撞击天王星所致\({^\text{[78]}}\)。在 1986 年旅行者 2 号探测器飞掠天王星时，其南极几乎正对太阳\({^\text{[79][80]}}\)。

自赫歇尔发现天卫三（Titania）与天卫四（Oberon）以来，天文学家便推测天王星有异常的自转轴，拉普拉斯在 1805 年计算了这些卫星轨道相对于行星赤道平面的倾斜度\({^\text{[81]}}\)。
\begin{figure}[ht]
\centering
\includegraphics[width=8cm]{./figures/733f8b68f8a5af74.png}
\caption{} \label{fig_Uranus_7}
\end{figure}
\subsubsection{从地球的可见性}
\begin{figure}[ht]
\centering
\includegraphics[width=8cm]{./figures/91965dcbef8408da.png}
\caption{“通过业余望远镜观测到的天王星，它出现在 2022 年 11 月月食期间月球掩星之后不久。”} \label{fig_Uranus_8}
\end{figure}
“天王星的平均视星等为 5.68，标准差为 0.17，而其亮度极值为 5.38 和 6.03\(^\text{[20]}\)。这一亮度范围接近肉眼可见的极限。亮度的大部分变化取决于从太阳照亮并从地球观测到的行星纬度\(^\text{[83]}\)。它的视直径介于 3.4 至 3.7 角秒之间，而土星为 16 至 20 角秒，木星为 32 至 45 角秒\(^\text{[84]}\)。在冲日时，天王星在暗夜条件下可被肉眼看到，即使在城市环境中也可以通过双筒望远镜轻松观测到\(^\text{[7]}\)。在口径 15 至 23 厘米的较大业余望远镜上，天王星呈现为一颗浅青色圆盘，并显示明显的边缘变暗。使用 25 厘米或更大口径的望远镜，可观测到云带结构，以及部分较大的卫星，如天卫三（Titania）和天卫四（Oberon）\(^\text{[85]}\)。
\subsubsection{内部结构}
\begin{figure}[ht]
\centering
\includegraphics[width=8cm]{./figures/a3223e5217a60cbf.png}
\caption{地球与天王星的大小比较} \label{fig_Uranus_9}
\end{figure}
天王星的质量约为地球的 14.5 倍，使其成为巨行星中质量最小者。其直径略大于海王星，约为地球的四倍。其由此产生的密度为 1.27 g/cm³，使天王星成为仅次于土星的第二低密度行星\(^\text{[11][12]}\)。这一数值表明它主要由各种“冰”组成，例如水、氨和甲烷\(^\text{[18]}\)。天王星内部冰的总质量并不精确已知，因为不同模型会得出不同数值；其必须介于 9.3 到 13.5 个地球质量之间\(^\text{[18][86]}\)。氢和氦只占总量的一小部分，约为 0.5 至 1.5 个地球质量\(^\text{[18]}\)。剩余的非冰物质（0.5 到 3.7 个地球质量）由岩石物质构成\(^\text{[18]}\)。

天王星的标准结构模型认为它由三层组成：中心为岩质（硅酸盐/铁–镍）核心，中间为冰质地幔，外层为气态的氢/氦包层\(^\text{[18][87]}\)。核心相对较小，质量仅为 0.55 个地球质量，半径不到行星的 20\%；地幔构成其主体，约为 13.4 个地球质量；上层大气相对不显著，质量约 0.5 个地球质量，延伸占天王星半径的最后 20\%\(^\text{[18][87]}\)。天王星核心的密度约为 9 g/cm³，中心压力约为 800 万巴（800 GPa），温度约为 5000 K\(^\text{[86][87]}\)。冰质地幔实际上并非传统意义的固态冰，而是由水、氨及其他挥发物组成的高温高密度流体\(^\text{[18][87]}\)。这种具有高电导率的流体有时被称为“水–氨海洋”\(^\text{[88]}\)。
\begin{figure}[ht]
\centering
\includegraphics[width=8cm]{./figures/9f28c21355c1f0ee.png}
\caption{天王星内部结构示意图，列出了各层的组成。} \label{fig_Uranus_10}
\end{figure}
在天王星深处的极端压力和温度条件下，甲烷分子可能会被分解，其中的碳原子会凝结成钻石晶体，并像冰雹一样穿过地幔降落\({^\text{[89][90]}}\)。这一现象类似于科学家推测可能在木星、土星和海王星上存在的“钻石雨”\({^\text{[91][92]}}\)。劳伦斯利弗莫尔国家实验室的超高压实验表明，在地幔底部可能存在一层金属液态碳的“海洋”，其中也许漂浮着固体“钻石冰山”\({^\text{[93][94][95]}}\)。

天王星和海王星的总体组成与木星和土星不同：冰的含量远多于气体，因此被单独归类为“冰巨行星”。在内部可能存在一层离子态水，其中水分子被分解为氢离子与氧离子的混合物；更深处则可能存在“超离子水”，在这种状态下氧原子形成晶格，而氢离子可以在晶格中自由移动\({^\text{[96]}}\)。

尽管上述结构模型被广泛采用，但并非唯一；其他模型同样能与观测数据相符。例如，如果大量氢和岩质物质混入冰状地幔，那么内部冰的总质量将减少，而相应的岩石和氢的总质量将增加。目前可获得的数据不足以科学地判定哪一种模型是正确的\({^\text{[86]}}\)。天王星内部为流体结构，因此并无固体表面；其气态大气会逐渐过渡到内部的液态层\({^\text{[18]}}\)。为了方便，在大气压等于 1 bar（100 kPa）的位置，约定一个旋转扁球体作为“表面”。该“表面”的赤道半径与极半径分别为 25,559 ± 4 km（15,881.6 ± 2.5 mi）和 24,973 ± 20 km（15,518 ± 12 mi）\({^\text{[11]}}\)。本条目中所有高度皆以此为零点。

\textbf{Internal heat} 天王星的内部热量显著低于其他巨行星；在天文学上，它具有很低的热通量\({^\text{[24][97]}}\)。为何天王星的内部温度如此之低仍然无法解释。与天王星在大小和成分上几乎相同的海王星向太空辐射的能量是其从太阳接收能量的 2.61 倍\({^\text{[24]}}\)，但天王星几乎不向外辐射多余热量。天王星在远红外（即热量）波段辐射的总功率为其大气吸收太阳能量的 1.06±0.08 倍\({^\text{[19][16]}}\)。天王星的热通量仅为 0.042±0.047 W/m\(^2\)，甚至低于地球约 0.075 W/m\(^2\) 的内部热通量\({^\text{[16]}}\)。在天王星对流层顶记录到的最低温度为 49 K（−224.2 °C；−371.5 °F），使其成为太阳系中最寒冷的行星\({^\text{[19][16]}}\)。

为了解释这种差异，一种假说认为，被认为导致天王星自转轴倾斜的地球大小的撞击体使天王星失去了大量原始热量，从而导致其核心温度显著降低\({^\text{[98]}}\)。另一种假说认为，天王星上层内部存在某种屏障，阻止核心热量向外传输\({^\text{[18]}}\)。例如，对不同成分层之间可能发生的对流过程的抑制可能会阻碍热量向上运输；\({^\text{[19][16]}}\) 双扩散对流可能是一种限制因素\({^\text{[18]}}\)。

在 2021 年的一项研究中，通过压缩含有橄榄石和铁方镁石等矿物的水，模拟了冰巨行星内部的条件，结果显示镁元素可以大量溶解在天王星和海王星的液态内部。如果天王星的镁含量高于海王星，它可能形成一层热绝缘层，从而潜在地解释该行星的低温\({^\text{[99]}}\)。

\textbf{Atmosphere}（大气） 尽管天王星内部不存在明确的固体表面，但天王星气态外壳中能够被遥感探测到的最外层区域被称为其大气\({^\text{[19]}}\)。遥感能力可深入到 1 bar（100 kPa）层以下约 300 km，对应压力约 100 bar（10 MPa）和温度约 320 K（47 °C；116 °F）\({^\text{[100]}}\)。稀薄的热层从名义上的 1 bar 表面向外延伸，覆盖超过两个行星半径\({^\text{[101]}}\)。天王星大气可分为三层：对流层，位于 −300 至 50 km（−186 至 31 mi）之间，压力从 100 bar 到 0.1 bar（10 MPa 到 10 kPa）；平流层，位于 50 至 4,000 km（31 至 2,485 mi）之间，压力从 0.1 bar 到 \(10^{-10}\) bar（10 kPa 到 10 μPa）；以及从 4,000 km 一直延伸到距表面约 50,000 km 的热层\({^\text{[19]}}\)。天王星不存在中间层（中间层即中气层）。

\textbf{Composition}（成分）



